\begin{document}
\begin{titlepage}
\centering
\includegraphics[width=0.8\textwidth]{NKU.png}\\[0.8cm]
{\Huge\kaishu\bfseries 南\quad 开\quad 大\quad 学\par}
\vspace{0.9cm}
{\huge\kaishu\bfseries 计\ 算\ 机\ 学\ 院\par}
\vspace{0.7cm}
{\Large\bfseries 操作系统第二次实验报告\par}
\vspace{0.8cm}
\HRule\par\vspace{0.7cm}
{\LARGE\bfseries 物理内存与页表\\[0.3cm]分配算法及 QEMU 调试\par}
\vspace{0.6cm}\HRule\par\vspace{1.2cm}
{\large\kaishu
\begin{tabular}{rl}
组长：&李珉熠\quad 2413982\\[0.5cm]
组员：&吴宇石\quad 2411080\\[0.5cm]
组员：&赵亚鑫\quad 2413715
\end{tabular}\par}
\vfill
{\Large 2026年10月23日\par}
\end{titlepage}
\pagenumbering{Roman}
\renewcommand{\abstractname}{\kaishu\sihao\textbf{摘要}}
\begin{abstract}
本实验在最小可执行内核基础上，完成物理内存探测与页描述符初始化、最佳适应页分配器、伙伴页分配器和简化对象分配器，并使用实际运行检查与课程评分验证结果。默认物理内存初始化及最佳适应分配评分为20/20；伙伴与对象分配另有独立测试。在分支任务中，我们构建独立带调试信息的QEMU，以客体GDB观察页表和控制寄存器，以宿主GDB跟踪取指及数据读取的地址翻译。实测启动映射采用顶层1GiB叶子，页表查询无需走满三级。报告区分课程提供的框架、实际修改、AI辅助过程与日志能够支持的结论。
\par\noindent\textbf{关键词：}物理内存；最佳适应；伙伴系统；对象分配；Sv39；QEMU调试
\end{abstract}
\clearpage
\tableofcontents
\clearpage
\pagenumbering{arabic}
\watermark
